\documentclass[11pt,a4paper]{article}

\usepackage[margin=1in]{geometry}
\usepackage{amsmath,amssymb,amsthm}
\usepackage{mathtools}
\usepackage{enumerate}
\usepackage{hyperref}
\usepackage{cleveref}
\usepackage[ruled,vlined,linesnumbered]{algorithm2e}
\usepackage{booktabs}

\newtheorem{theorem}{Theorem}[section]
\newtheorem{lemma}[theorem]{Lemma}
\newtheorem{proposition}[theorem]{Proposition}
\newtheorem{corollary}[theorem]{Corollary}
\newtheorem{definition}[theorem]{Definition}
\newtheorem{remark}[theorem]{Remark}

\DeclareMathOperator{\proj}{proj}
\DeclareMathOperator{\AM}{AM}
\DeclareMathOperator{\HM}{HM}
\DeclareMathOperator{\diag}{diag}
\DeclareMathOperator{\tr}{tr}
\DeclareMathOperator*{\argmin}{arg\,min}
\newcommand{\R}{\mathbb{R}}
\newcommand{\E}{\mathbb{E}}
\newcommand{\Prob}{\mathbb{P}}
\newcommand{\cF}{\mathcal{F}}
\newcommand{\cG}{\mathcal{G}}
\newcommand{\cS}{\mathcal{S}}
\newcommand{\cC}{\mathcal{C}}
\newcommand{\cD}{\mathcal{D}}
\newcommand{\ip}[2]{\langle #1, #2 \rangle}
\newcommand{\norm}[1]{\|#1\|}

\title{Evidence-Weighted Federated Averaging:\\
An Application of the $\Omega$-Framework\\to Federated Learning}
\author{Eugene Shcherbinin\\
\small BSc Mathematics, London School of Economics\\
\small Candidate 007}
\date{March 2026}

\begin{document}
\maketitle

\begin{abstract}
We apply the $\Omega$-framework (Evidence-Weighted Multi-Agent Policy Gradient)
to federated learning (FL). The mapping is precise: FL clients are agents,
data heterogeneity across clients is evidence variance $V_i$, gradient
compression is the lossy communication channel, and client selection is
coalition formation. Replacing FedAvg's uniform aggregation with the evidence
weight $w_i = V_{\min}/V_i$ yields \emph{Evidence-Weighted Federated Averaging}
(EW-FedAvg), which achieves a variance improvement factor of
$\HM(V)/\AM(V) \leq 1$ over standard FedAvg. The self-knowledge bottleneck
from the cooperative PG framework transfers directly: clients with highly
non-IID data cannot improve their contribution to the global model through
bandwidth alone --- they need more \emph{local data}. We prove communication-adjusted
convergence rates, connect sparsified gradient communication to the blessing
of dimensionality, and show that evidence-weighted client selection subsumes
and unifies the robustness mechanisms of SCAFFOLD and FedProx.
\end{abstract}

\tableofcontents

%% ========================================================================
\section{Introduction}
\label{sec:intro}
%% ========================================================================

Federated learning~\cite{mcmahan2017} trains a shared model across $K$
clients without centralizing data. The canonical algorithm, FedAvg,
aggregates client updates using weights proportional to local dataset size
$n_i / n$, where $n = \sum_i n_i$. This weighting is optimal when client
data distributions are identical (IID). In practice, client data is
heterogeneous (non-IID), and uniform-style aggregation produces a biased
global model that converges slowly or diverges~\cite{karimireddy2020,li2020fedprox}.

The $\Omega$-framework~\cite{ewpg2026} was developed for multi-agent
reinforcement learning, where agents with heterogeneous observation quality
coordinate through lossy channels. Its core mechanism --- the evidence weight
$w_i = V_{\min}/V_i$ measuring each agent's information quality --- transfers
to federated learning with minimal modification. In this paper, we make this
transfer precise and derive its consequences.

\textbf{Contributions.}
\begin{enumerate}[(1)]
  \item A formal isomorphism between the $\Omega$-framework and federated
    learning (\Cref{sec:isomorphism}).
  \item Evidence-Weighted FedAvg with a $\HM(V)/\AM(V)$ variance improvement
    theorem (\Cref{sec:ew-fedavg}).
  \item Communication efficiency analysis via the self-knowledge bottleneck,
    with connections to sparse recovery (\Cref{sec:communication}).
  \item Client selection as coalition formation with evidence-weighted
    rationality criteria (\Cref{sec:coalition}).
  \item Non-IID robustness analysis unifying SCAFFOLD and FedProx under the
    evidence-weighting lens (\Cref{sec:noniid}).
\end{enumerate}


%% ========================================================================
\section{The Isomorphism}
\label{sec:isomorphism}
%% ========================================================================

\subsection{Setup: Federated Learning}

Consider $K$ clients, each holding a local dataset $\cD_i$ drawn from a
(possibly distinct) distribution $P_i$ over input-label pairs $(x, y)$.
The global objective is:
\begin{equation}\label{eq:fl-objective}
  \min_{\theta \in \R^d}\; F(\theta) := \sum_{i=1}^{K} p_i\, F_i(\theta),
  \qquad F_i(\theta) := \E_{(x,y) \sim P_i}[\ell(\theta; x, y)],
\end{equation}
where $p_i \geq 0$, $\sum_i p_i = 1$ are aggregation weights (typically
$p_i = n_i / n$), and $\ell$ is the loss function.

At each communication round $t$, FedAvg~\cite{mcmahan2017} proceeds:
\begin{enumerate}[(i)]
  \item Server broadcasts current model $\theta^t$ to selected clients.
  \item Each selected client $i$ runs $\tau$ steps of SGD on $F_i$,
    producing a local update $\Delta_i^t = \theta_i^t - \theta^t$.
  \item Server aggregates: $\theta^{t+1} = \theta^t + \sum_{i \in S_t} p_i \Delta_i^t$.
\end{enumerate}

\subsection{The $\Omega$--FL Correspondence}

\begin{table}[h]
\centering
\caption{Precise mapping between the $\Omega$-framework and federated learning.}
\label{tab:isomorphism}
\begin{tabular}{@{}lll@{}}
\toprule
\textbf{$\Omega$-Framework} & \textbf{Federated Learning} & \textbf{Formal Correspondence} \\
\midrule
Agent $i \in \{1,\ldots,N\}$ & Client $i \in \{1,\ldots,K\}$ & $N = K$ \\[3pt]
Policy $\pi_i \in \Pi_i$ & Local model $\theta_i \in \R^d$ & $\Pi_i = \R^d$ \\[3pt]
Joint policy $\pi \in \Pi$ & Global model $\theta \in \R^d$ & $\theta = \sum_i p_i \theta_i$ \\[3pt]
Policy gradient $\hat{v}_i$ & Local gradient $g_i = \nabla F_i(\theta)$ & $\hat{v}_i \leftrightarrow g_i$ \\[3pt]
Evidence variance $V_i$ & Gradient variance $\sigma_i^2$ & $V_i = \sigma_i^2$ \\[3pt]
 & ~~$= \E[\norm{g_i - \nabla F(\theta)}^2]$ & \\[3pt]
Evidence weight $w_i = V_{\min}/V_i$ & Client quality weight & $w_i = \sigma_{\min}^2 / \sigma_i^2$ \\[3pt]
Lossy comm.\ channel & Gradient compression & Quantization, sparsification \\[3pt]
Self-knowledge $\hat{\pi}_i$ & Client's local model & Drift $\norm{\theta_i - \theta}$ \\[3pt]
Self-knowledge loss $L_{\text{self}}^{(i)}$ & Client drift & $\E[\norm{\theta_i^t - \theta^t}^2]$ \\[3pt]
Coalition $S \subseteq \{1,\ldots,N\}$ & Selected client set $S_t$ & $S = S_t$ \\[3pt]
Nash equilibrium $\pi^*$ & Global optimum $\theta^*$ & $\nabla F(\theta^*) = 0$ \\[3pt]
Lyapunov $D(\pi) = \frac{1}{2}\norm{\pi - \pi^*}^2$
  & Suboptimality $F(\theta) - F(\theta^*)$ & Convergence measure \\
\bottomrule
\end{tabular}
\end{table}

\begin{remark}[Evidence variance = gradient heterogeneity]
The evidence variance $V_i$ in the $\Omega$-framework measures the
noise-to-signal ratio of agent~$i$'s gradient estimate. In FL, the
natural counterpart is the \emph{gradient divergence}:
\begin{equation}\label{eq:grad-div}
  \sigma_i^2 := \E\bigl[\norm{\nabla F_i(\theta) - \nabla F(\theta)}^2\bigr],
\end{equation}
which is driven by two sources: (a)~stochastic noise from mini-batch
sampling, and (b)~\emph{distribution heterogeneity} $P_i \neq P_j$.
Under IID data, $\sigma_i^2$ reduces to mini-batch variance; under
non-IID data, it is dominated by the distribution shift. This is
precisely the structure the $\Omega$-framework was built to handle.
\end{remark}

\begin{definition}[Gradient heterogeneity measures]\label{def:heterogeneity}
For FL with $K$ clients:
\begin{align}
  \sigma_i^2 &:= \E\bigl[\norm{\nabla F_i(\theta) - \nabla F(\theta)}^2\bigr]
    &&\text{(client-level variance)}, \label{eq:sigma-i}\\
  \bar{\sigma}^2 &:= \sum_{i=1}^{K} p_i\, \sigma_i^2
    &&\text{(average variance = $\AM(\sigma^2)$)}, \label{eq:sigma-bar}\\
  \sigma_{\min}^2 &:= \min_{i} \sigma_i^2
    &&\text{(best-informed client)}. \label{eq:sigma-min}
\end{align}
\end{definition}


%% ========================================================================
\section{Evidence-Weighted Federated Averaging}
\label{sec:ew-fedavg}
%% ========================================================================

\subsection{The Aggregation Rule}

Standard FedAvg aggregates with weights $p_i = n_i / n$ (proportional to
data size). We replace this with the evidence-weighted aggregation:

\begin{definition}[Evidence-Weighted FedAvg]\label{def:ew-fedavg}
At communication round $t$, the server computes:
\begin{equation}\label{eq:ew-agg}
  \theta^{t+1} = \theta^t + \sum_{i \in S_t} \tilde{w}_i\, \Delta_i^t,
  \qquad \tilde{w}_i = \frac{w_i}{\sum_{j \in S_t} w_j},
  \qquad w_i = \frac{\sigma_{\min}^2}{\sigma_i^2}.
\end{equation}
\end{definition}

\noindent Clients with lower gradient variance (higher evidence quality)
receive higher aggregation weight. The normalization $\sum_i \tilde{w}_i = 1$
preserves the averaging property. When all $\sigma_i^2$ are equal, $w_i = 1$
for all~$i$ and EW-FedAvg reduces to uniform averaging.

\subsection{Variance Improvement}

\begin{theorem}[EW-FedAvg variance improvement]\label{thm:ew-variance}
Let $g_i^t = \nabla F_i(\theta^t) + \xi_i^t$ be the stochastic gradient of
client~$i$, with $\E[\xi_i^t] = 0$ and
$\E[\norm{g_i^t - \nabla F(\theta^t)}^2] = \sigma_i^2$. Then the aggregated
gradient under EW-FedAvg satisfies:
\begin{equation}\label{eq:ew-var-bound}
  \E\Bigl[\Bigl\|\sum_{i=1}^{K} \tilde{w}_i\, g_i^t - \nabla F(\theta^t)\Bigr\|^2\Bigr]
  \;\leq\; \frac{\HM(\sigma^2)}{\AM(\sigma^2)} \cdot \bar{\sigma}^2
  \;\leq\; \bar{\sigma}^2,
\end{equation}
where $\HM(\sigma^2) = K / \bigl(\sum_{i=1}^K 1/\sigma_i^2\bigr)$ and
$\AM(\sigma^2) = (1/K)\sum_{i=1}^K \sigma_i^2$. The ratio
$\HM(\sigma^2)/\AM(\sigma^2) \leq 1$ with equality iff all $\sigma_i^2$ are
identical.
\end{theorem}

\begin{proof}
The aggregated gradient error is:
\begin{align}
  \E\Bigl[\Bigl\|\sum_i \tilde{w}_i g_i - \nabla F\Bigr\|^2\Bigr]
  &= \sum_i \tilde{w}_i^2 \,\sigma_i^2
  \label{eq:ew-step1}
\end{align}
where the cross terms vanish since gradient errors across clients are
independent conditional on $\theta^t$ (clients sample from independent
local datasets). The weights are $\tilde{w}_i = w_i / W$ with
$W = \sum_j w_j = \sigma_{\min}^2 \sum_j 1/\sigma_j^2$, so:
\begin{equation}
  \tilde{w}_i^2 \,\sigma_i^2
  = \frac{(\sigma_{\min}^2)^2}{\sigma_i^4} \cdot \frac{\sigma_i^2}{W^2}
  = \frac{(\sigma_{\min}^2)^2}{\sigma_i^2 \, W^2}.
\end{equation}
Summing:
\begin{equation}
  \sum_i \tilde{w}_i^2 \sigma_i^2
  = \frac{(\sigma_{\min}^2)^2}{W^2} \sum_i \frac{1}{\sigma_i^2}
  = \frac{(\sigma_{\min}^2)^2}{W^2} \cdot \frac{W}{\sigma_{\min}^2}
  = \frac{\sigma_{\min}^2}{W}
  = \frac{1}{\sum_i 1/\sigma_i^2}.
  \label{eq:ew-sum}
\end{equation}
This is $\HM(\sigma^2) / K$. Now, uniform averaging gives
$\bar{\sigma}^2 / K = \AM(\sigma^2) / K$ (each client weighted $1/K$,
contributing $\sigma_i^2 / K^2$, summing to $\AM(\sigma^2)/K$). Therefore:
\begin{equation}
  \frac{\text{EW-FedAvg variance}}{\text{FedAvg variance}}
  = \frac{\HM(\sigma^2) / K}{\AM(\sigma^2) / K}
  = \frac{\HM(\sigma^2)}{\AM(\sigma^2)} \leq 1,
\end{equation}
where the final inequality is the HM-AM inequality. Equality holds iff
all $\sigma_i^2$ are equal.
\end{proof}

\begin{remark}[Magnitude of improvement]
If $K = 2$ clients have $\sigma_1^2 = 1$ and $\sigma_2^2 = 100$, then:
\[
  \frac{\HM}{\AM} = \frac{2/(1 + 1/100)}{(1+100)/2}
  = \frac{200/101}{101/2} = \frac{400}{10201} \approx 0.039.
\]
The variance is reduced by a factor of $\sim\!25\times$. Highly heterogeneous
settings produce the largest gains --- precisely the regime where FedAvg
struggles most.
\end{remark}

\subsection{Convergence of EW-FedAvg}

\begin{theorem}[Convergence rate of EW-FedAvg]\label{thm:ew-convergence}
Assume $F$ is $L$-smooth, each $F_i$ is $\mu$-strongly convex, and
$\E[\norm{g_i - \nabla F_i}^2] \leq \sigma_{\text{sgd}}^2$ (mini-batch
noise). After $T$ communication rounds with $\tau$ local SGD steps each
and step size $\eta = 1 / (4L\tau)$:
\begin{equation}\label{eq:ew-conv-rate}
  \E[F(\theta^T)] - F^*
  \;\leq\; \underbrace{\left(1 - \frac{\mu}{4L}\right)^T}_{\text{linear contraction}}
  \cdot (F(\theta^0) - F^*)
  \;+\; \frac{8L}{\mu^2}\left(
    \frac{\sigma_{\text{sgd}}^2}{K\tau}
    + \tau\,\frac{\HM(\sigma^2)}{\AM(\sigma^2)}\,\bar{\sigma}^2
  \right).
\end{equation}
Compare FedAvg~\cite{karimireddy2020}, where the second floor term is
$\tau\,\bar{\sigma}^2$ (no $\HM/\AM$ factor).
\end{theorem}

\begin{proof}[Proof sketch]
Follow the standard FedAvg analysis of Karimireddy et al.~\cite{karimireddy2020},
replacing uniform aggregation weights $p_i = 1/K$ with evidence weights
$\tilde{w}_i$. The $L$-smoothness descent lemma gives:
\[
  \E[F(\theta^{t+1})] \leq F(\theta^t)
    - \frac{\eta\tau}{2}\norm{\nabla F(\theta^t)}^2
    + \frac{L\eta^2\tau}{2}\sigma_{\text{sgd}}^2
    + \frac{L\eta^2\tau^2}{2}\sum_i \tilde{w}_i^2\,\sigma_i^2.
\]
The final term is bounded by $\HM(\sigma^2)/(K) \cdot L\eta^2\tau^2 / 2$
via~\eqref{eq:ew-sum}. Telescoping and using strong convexity yields
\eqref{eq:ew-conv-rate}.
\end{proof}

\subsection{Algorithm}

\begin{algorithm}[H]
\caption{Evidence-Weighted Federated Averaging (EW-FedAvg)}
\label{alg:ew-fedavg}
\SetAlgoLined
\KwIn{Initial model $\theta^0$, $K$ clients, local epochs $\tau$,
  learning rate $\eta$, communication rounds $T$}
\KwOut{Trained global model $\theta^T$}

\tcp{Server estimates client variances (one-time or periodic)}
\For{$i = 1, \ldots, K$}{
  $\hat{\sigma}_i^2 \gets \textsc{EstimateVariance}(i, \theta^0)$\;
}
$\hat{\sigma}_{\min}^2 \gets \min_i \hat{\sigma}_i^2$\;

\BlankLine
\For{$t = 0, 1, \ldots, T-1$}{
  \tcp{Client selection (coalition formation)}
  $S_t \gets \textsc{SelectClients}(\{w_i\}_{i=1}^K)$\;
  Broadcast $\theta^t$ to all $i \in S_t$\;

  \BlankLine
  \tcp{Local training}
  \For(\tcp*[f]{in parallel}){each $i \in S_t$}{
    $\theta_i^{t,0} \gets \theta^t$\;
    \For{$k = 0, \ldots, \tau - 1$}{
      Sample mini-batch $\xi_i \sim \cD_i$\;
      $\theta_i^{t,k+1} \gets \theta_i^{t,k} - \eta\,\nabla F_i(\theta_i^{t,k}; \xi_i)$\;
    }
    $\Delta_i^t \gets \theta_i^{t,\tau} - \theta^t$\;
    \tcp{Optional: sparsify $\Delta_i^t$ (Section~\ref{sec:communication})}
  }

  \BlankLine
  \tcp{Evidence-weighted aggregation}
  $w_i \gets \hat{\sigma}_{\min}^2 / \hat{\sigma}_i^2$ for each $i \in S_t$\;
  $W \gets \sum_{i \in S_t} w_i$\;
  $\theta^{t+1} \gets \theta^t + \sum_{i \in S_t} \frac{w_i}{W}\,\Delta_i^t$\;
}
\Return{$\theta^T$}\;
\end{algorithm}


%% ========================================================================
\section{Communication Efficiency and the Self-Knowledge Bottleneck}
\label{sec:communication}
%% ========================================================================

\subsection{The Self-Knowledge Bottleneck in FL}

The cooperative PG framework established that an agent's ability to
communicate its policy is fundamentally limited by its self-knowledge.
In FL, the analogous statement is:

\begin{definition}[Client drift as self-knowledge loss]\label{def:drift}
After $\tau$ local SGD steps, client~$i$'s drift from the global model is:
\begin{equation}\label{eq:drift}
  \Delta_{\text{drift}}^{(i)} := \E\bigl[\norm{\theta_i^{t,\tau} - \theta^t}^2\bigr].
\end{equation}
This is the FL analogue of the self-knowledge loss $L_{\text{self}}^{(i)}$:
the local model $\theta_i$ is client~$i$'s ``self-model'' of its
contribution to the global objective.
\end{definition}

\begin{proposition}[Drift--heterogeneity relationship]\label{prop:drift-bound}
Under $L$-smooth local objectives with step size $\eta \leq 1/L$:
\begin{equation}\label{eq:drift-var}
  \Delta_{\text{drift}}^{(i)} \;\leq\; \tau^2 \eta^2\,\sigma_i^2 + \tau\eta^2\,\sigma_{\text{sgd}}^2.
\end{equation}
The first term is driven by distribution heterogeneity; the second by
stochastic noise.
\end{proposition}

\begin{proof}
By $L$-smoothness and the SGD recursion:
\begin{align}
  \theta_i^{t,k+1} - \theta^t
  &= \theta_i^{t,k} - \theta^t - \eta\,\nabla F_i(\theta_i^{t,k}; \xi_i^k).
\end{align}
Taking expectations and using $\E[\nabla F_i(\theta_i^{t,k}; \xi)] = \nabla F_i(\theta_i^{t,k})$,
then bounding $\norm{\nabla F_i(\theta^t) - \nabla F(\theta^t)} \leq \sigma_i$
and applying the standard SGD drift analysis~\cite{karimireddy2020}
yields~\eqref{eq:drift-var}.
\end{proof}

\begin{corollary}[Self-knowledge bottleneck in FL]\label{cor:fl-bottleneck}
The useful information in client~$i$'s update $\Delta_i^t$ is bounded by:
\begin{equation}
  I(\theta^*;\, \Delta_i^t) \;\leq\; I(\theta^*;\, \cD_i).
\end{equation}
Even with \emph{infinite bandwidth} (no compression, no quantization),
client~$i$'s contribution to the global model is limited by the information
content of its local data about the global optimum.

For a highly non-IID client (large $\sigma_i^2$), the local data contains
limited information about $\theta^*$. No amount of communication bandwidth
compensates --- the client needs more \emph{data}, not more \emph{bandwidth}.
\end{corollary}

\subsection{Communication-Adjusted Convergence}

When clients communicate compressed updates $\hat{\Delta}_i^t$ instead of
exact updates $\Delta_i^t$, a channel distortion term enters:
\begin{equation}\label{eq:channel-loss}
  L_{\text{channel}}^{(i)} := \E\bigl[\norm{\Delta_i^t - \hat{\Delta}_i^t}^2\bigr].
\end{equation}

\begin{theorem}[Communication-adjusted convergence]\label{thm:comm-conv}
Under the conditions of \Cref{thm:ew-convergence}, with compressed
communication satisfying $L_{\text{channel}}^{(i)} \leq \delta_i^2\norm{\Delta_i^t}^2$
for compression ratio $\delta_i \in [0,1]$, EW-FedAvg converges at rate:
\begin{equation}\label{eq:comm-conv}
  \E[F(\theta^T)] - F^*
  \;\leq\; \left(1 - \frac{\mu}{4L}\right)^T (F(\theta^0) - F^*)
  + \frac{8L}{\mu^2}\left(
    \frac{\sigma_{\text{sgd}}^2}{K\tau}
    + \tau\,\frac{\HM(\sigma^2)}{\AM(\sigma^2)}\,\bar{\sigma}^2
    + \sum_i \tilde{w}_i^2\,\delta_i^2\,\Delta_{\text{drift}}^{(i)}
  \right).
\end{equation}
The third floor term decomposes via~\eqref{eq:drift-var} into a
self-knowledge component ($\tau^2\eta^2\sigma_i^2$) and a stochastic
component ($\tau\eta^2\sigma_{\text{sgd}}^2$), both modulated by the
compression ratio~$\delta_i$ and the evidence weight~$\tilde{w}_i$.
\end{theorem}

\begin{proof}[Proof sketch]
The compressed update adds noise:
$\hat{\Delta}_i^t = \Delta_i^t + \epsilon_i^t$ with
$\E[\norm{\epsilon_i^t}^2] \leq \delta_i^2\norm{\Delta_i^t}^2$.
This enters the descent lemma as an additional variance term:
\[
  \E\bigl[\norm{\textstyle\sum_i \tilde{w}_i \epsilon_i^t}^2\bigr]
  = \sum_i \tilde{w}_i^2 \E[\norm{\epsilon_i^t}^2]
  \leq \sum_i \tilde{w}_i^2 \delta_i^2 \Delta_{\text{drift}}^{(i)}.
\]
Combining with the EW-FedAvg analysis of \Cref{thm:ew-convergence}
gives~\eqref{eq:comm-conv}.
\end{proof}

\begin{remark}[Evidence-aware compression]
\Cref{thm:comm-conv} suggests an optimal compression strategy: allocate
more bandwidth to high-evidence clients (low $\sigma_i^2$, high $\tilde{w}_i$)
since their updates carry more weight. Formally, to minimize the
communication floor $\sum_i \tilde{w}_i^2 \delta_i^2 \Delta_{\text{drift}}^{(i)}$
subject to a total bandwidth constraint $\sum_i R_i \leq R_{\text{total}}$,
the optimal allocation is water-filling:
\begin{equation}
  R_i^* \;\propto\; \tilde{w}_i \sqrt{\Delta_{\text{drift}}^{(i)}}.
\end{equation}
High-evidence clients get more bandwidth; low-evidence clients (whose
updates are downweighted anyway) can be aggressively compressed.
\end{remark}

\subsection{Sparsified Gradients and the Blessing of Dimensionality}

The blessing of dimensionality from the $\Omega$-framework~\cite{blessing2026}
applies directly to FL gradient communication. When the global model has
$d$ parameters but the effective gradient at each client is $k$-sparse
(concentrates on $k \ll d$ coordinates), compressed sensing guarantees
apply:

\begin{proposition}[Sparse gradient recovery]\label{prop:sparse-recovery}
If client~$i$'s true gradient $\nabla F_i(\theta)$ is $k_i$-sparse (has
at most $k_i$ non-negligible coordinates), then top-$k_i$ sparsification
achieves:
\begin{equation}\label{eq:sparse-rate}
  L_{\text{channel}}^{(i)} \;\leq\; \frac{\sigma_{\text{sgd}}^2}{n_i}
    \cdot \frac{k_i \log(d/k_i)}{d},
\end{equation}
requiring only $O(k_i \log(d/k_i))$ communicated coordinates instead of $d$.
\end{proposition}

The key insight: non-IID clients tend to have \emph{sparser} gradients.
A client whose data covers only a narrow slice of the input space has
gradients concentrated on the model parameters relevant to that slice.
This means the clients that need the most compression (high $\sigma_i^2$,
heavily non-IID) are also the most compressible.

\begin{corollary}[Sparsity--heterogeneity duality]
Gradient sparsity $k_i$ and evidence variance $\sigma_i^2$ are positively
correlated across clients: $k_i = O(1/\sigma_i^2)$ when heterogeneity
is driven by data support mismatch. The evidence weight $w_i$ thus
simultaneously controls aggregation quality and compression efficiency.
\end{corollary}


%% ========================================================================
\section{Client Selection as Coalition Formation}
\label{sec:coalition}
%% ========================================================================

\subsection{Coalition Rationality in FL}

In the $\Omega$-framework, coalition~$S$ is rational if its members
collectively outperform their sum of individual values. In FL, the
analogous question is: which subset $S_t$ of clients should the server
select at round~$t$?

\begin{definition}[Marginal contribution of client $i$]\label{def:marginal}
The marginal contribution of client~$i$ to coalition $S_t$ is:
\begin{equation}\label{eq:marginal}
  \Delta V_i(S_t) := \E[F(\theta_{S_t}^{t+1})] - \E[F(\theta_{S_t \setminus \{i\}}^{t+1})],
\end{equation}
where $\theta_S^{t+1}$ denotes the model obtained by aggregating updates
from clients in~$S$ only.
\end{definition}

\begin{definition}[Communication-adjusted contribution]\label{def:comm-adj}
Accounting for communication cost $c_i$ (bandwidth, latency, compute),
client~$i$ should be included in $S_t$ iff:
\begin{equation}\label{eq:rational-inclusion}
  |\Delta V_i(S_t)| > c_i.
\end{equation}
\end{definition}

\subsection{Evidence-Weighted Selection}

\begin{theorem}[Evidence weight as selection criterion]\label{thm:selection}
Under $L$-smoothness and the conditions of \Cref{thm:ew-variance}, the
marginal contribution of client~$i$ satisfies:
\begin{equation}\label{eq:marginal-bound}
  |\Delta V_i(S_t)| \;\leq\; \frac{L\eta\tau}{W^2}\,w_i
  \left(\norm{\nabla F(\theta^t)}^2 + \sigma_i^2\right),
\end{equation}
where $W = \sum_{j \in S_t} w_j$. The contribution is proportional to
the evidence weight~$w_i$: low-evidence clients (high $\sigma_i^2$, low
$w_i$) contribute less to the aggregate and are naturally deprioritized.
\end{theorem}

\begin{proof}[Proof sketch]
The EW-FedAvg aggregated update with client~$i$ included vs.\ excluded
differs by the weighted increment $(\tilde{w}_i / W) \Delta_i^t$. By
$L$-smoothness:
\[
  |F(\theta_{S_t}^{t+1}) - F(\theta_{S_t \setminus i}^{t+1})|
  \leq L\norm{\theta_{S_t}^{t+1} - \theta_{S_t \setminus i}^{t+1}}
  = L \cdot \frac{w_i}{W} \norm{\Delta_i^t}.
\]
Bounding $\E[\norm{\Delta_i^t}] \leq \eta\tau(\norm{\nabla F(\theta^t)} + \sigma_i)$
and collecting terms gives~\eqref{eq:marginal-bound}.
\end{proof}

\begin{corollary}[Automatic client filtering]\label{cor:auto-filter}
Evidence weighting provides a soft client selection mechanism without
explicit selection: even if all $K$ clients participate, the evidence
weights $\tilde{w}_i = w_i / W$ automatically suppress the contribution
of low-quality clients. This is strictly better than hard selection
(which discards information) while achieving the same robustness.
\end{corollary}

\subsection{Cluster-Based Coalitions}

When clients can be grouped into clusters $\cC_1, \ldots, \cC_M$ with
similar distributions ($\sigma_i^2 \approx \sigma_j^2$ for $i, j$ in the
same cluster), a two-level aggregation emerges naturally:

\begin{enumerate}[(i)]
  \item \textbf{Intra-cluster}: average within each cluster (clients are
    approximately IID within a cluster, so uniform weighting is near-optimal).
  \item \textbf{Inter-cluster}: aggregate cluster-level updates using
    evidence weights $w_m = \sigma_{\min}^2 / \bar{\sigma}_m^2$, where
    $\bar{\sigma}_m^2$ is cluster~$m$'s average variance.
\end{enumerate}

This corresponds precisely to the hierarchical coalition structure in
the $\Omega$-framework: intra-cluster cooperation (local coalitions)
followed by inter-cluster weighted aggregation (global coalition).


%% ========================================================================
\section{Non-IID Robustness}
\label{sec:noniid}
%% ========================================================================

\subsection{The Non-IID Problem}

The central challenge in FL is data heterogeneity. When $P_i \neq P_j$,
the local objectives $F_i$ have different optima $\theta_i^* \neq \theta_j^*$,
and local SGD steps drive clients away from the global optimum~$\theta^*$.
This ``client drift'' is the dominant source of convergence degradation.

Existing approaches address this through:
\begin{itemize}
  \item \textbf{SCAFFOLD}~\cite{karimireddy2020}: control variates
    $c_i$ that correct for the difference $\nabla F_i - \nabla F$ at each
    client.
  \item \textbf{FedProx}~\cite{li2020fedprox}: a proximal regularizer
    $(\mu_{\text{prox}}/2)\norm{\theta - \theta^t}^2$ added to each local
    objective, penalizing drift from the global model.
\end{itemize}

\subsection{Evidence Weighting as Unified Robustness Mechanism}

\begin{proposition}[EW-FedAvg subsumes FedProx's effect]\label{prop:fedprox}
FedProx's proximal term reduces client drift:
$\Delta_{\text{drift}}^{(i),\text{FedProx}} \leq \Delta_{\text{drift}}^{(i)} / (1 + \mu_{\text{prox}}/L)^2$.
EW-FedAvg achieves an analogous effect through aggregation: the contribution
of drifted clients to the global update is suppressed by $w_i = \sigma_{\min}^2/\sigma_i^2$.

Formally, the effective drift contribution under EW-FedAvg is:
\begin{equation}\label{eq:ew-vs-fedprox}
  \sum_i \tilde{w}_i \,\Delta_{\text{drift}}^{(i)}
  \;\leq\; \frac{\HM(\sigma^2)}{\AM(\sigma^2)} \cdot \bar{\Delta}_{\text{drift}},
\end{equation}
where $\bar{\Delta}_{\text{drift}} = \sum_i (1/K)\,\Delta_{\text{drift}}^{(i)}$.
\end{proposition}

\begin{proof}
Since $\Delta_{\text{drift}}^{(i)} \leq C\,\sigma_i^2$ for some constant
$C$ depending on $\eta, \tau$ (\Cref{prop:drift-bound}):
\begin{equation}
  \sum_i \tilde{w}_i\,\Delta_{\text{drift}}^{(i)}
  \leq C \sum_i \tilde{w}_i\,\sigma_i^2
  = C \sum_i \frac{w_i}{W}\,\sigma_i^2
  = C \sum_i \frac{\sigma_{\min}^2}{W}
  = \frac{CK\sigma_{\min}^2}{W}
  = \frac{C \cdot \HM(\sigma^2)}{1}.
\end{equation}
Meanwhile, $\bar{\Delta}_{\text{drift}} \leq C \cdot \AM(\sigma^2)$,
giving the ratio~\eqref{eq:ew-vs-fedprox}.
\end{proof}

\begin{proposition}[EW-FedAvg complements SCAFFOLD]\label{prop:scaffold}
SCAFFOLD's control variates correct the bias: $\E[g_i - c_i] = \nabla F(\theta)$.
This reduces $\sigma_i^2$ to the residual variance after control-variate
correction:
\begin{equation}
  \sigma_i^{2,\text{SCAFFOLD}} = \E\bigl[\norm{g_i - c_i - \nabla F(\theta)}^2\bigr]
  \;\leq\; \sigma_{\text{sgd}}^2.
\end{equation}
When $\sigma_i^{2,\text{SCAFFOLD}} \approx \sigma_{\text{sgd}}^2$ for all
clients, evidence weights become approximately uniform: $w_i \approx 1$.
EW-FedAvg gracefully reduces to standard averaging when heterogeneity is
already controlled. The two mechanisms are complementary:
\begin{center}
\begin{tabular}{lll}
\toprule
\textbf{Mechanism} & \textbf{Target} & \textbf{Action} \\
\midrule
SCAFFOLD & Bias $\E[g_i] \neq \nabla F$ & Corrects via control variate \\
EW-FedAvg & Variance $\text{Var}(g_i)$ heterogeneity & Reweights via $w_i$ \\
FedProx & Client drift $\norm{\theta_i - \theta}$ & Penalizes via regularizer \\
\bottomrule
\end{tabular}
\end{center}
\end{proposition}

\subsection{Convergence Under Arbitrary Heterogeneity}

\begin{theorem}[Heterogeneity-adaptive convergence]\label{thm:noniid-conv}
Under the conditions of \Cref{thm:ew-convergence}, EW-FedAvg achieves:
\begin{equation}
  \E[F(\theta^T)] - F^* \;\leq\;
  \left(1 - \frac{\mu}{4L}\right)^T (F(\theta^0) - F^*)
  + \frac{8L\tau}{\mu^2}\,\frac{K}{\sum_{i=1}^K 1/\sigma_i^2}
  + \frac{8L}{\mu^2} \cdot \frac{\sigma_{\text{sgd}}^2}{K\tau}.
\end{equation}
In the homogeneous limit ($\sigma_i^2 = \sigma^2$ for all~$i$), this
recovers FedAvg's rate with floor $8L\tau\sigma^2/\mu^2$. In the
heterogeneous regime, the floor is determined by the \emph{harmonic mean}
of the variances --- dominated by the best clients rather than the worst.
\end{theorem}

\begin{remark}[Comparison with existing bounds]
FedAvg's convergence floor scales with $\AM(\sigma^2)$ --- the worst clients
dominate. FedProx scales with $\AM(\sigma^2)/(1 + \mu_{\text{prox}}/L)^2$
--- reduced but still dominated by worst clients. EW-FedAvg scales with
$\HM(\sigma^2)$ --- dominated by the \emph{best} clients. This is the
fundamental advantage: evidence weighting leverages information quality
differences rather than suffering from them.
\end{remark}


%% ========================================================================
\section{The Full Algorithm}
\label{sec:full-algorithm}
%% ========================================================================

Combining all components --- evidence-weighted aggregation, communication
efficiency, and adaptive client selection --- we obtain the full
Evidence-Weighted Federated Averaging algorithm.

\begin{algorithm}[H]
\caption{Full EW-FedAvg with Adaptive Selection and Compression}
\label{alg:full-ew-fedavg}
\SetAlgoLined
\KwIn{Initial model $\theta^0$, $K$ clients with datasets $\{\cD_i\}$,
  local epochs $\tau$, step size $\eta$, rounds $T$,
  bandwidth budget $R_{\text{total}}$, cost threshold $c_{\min}$}
\KwOut{Trained global model $\theta^T$}

\BlankLine
\tcp{Phase 0: Variance estimation (run periodically)}
\For{$i = 1, \ldots, K$}{
  Compute $\hat{\sigma}_i^2$ from pilot gradients at $\theta^0$\;
}
$\hat{\sigma}_{\min}^2 \gets \min_i \hat{\sigma}_i^2$\;
$w_i \gets \hat{\sigma}_{\min}^2 / \hat{\sigma}_i^2$ for each $i$\;

\BlankLine
\For{$t = 0, 1, \ldots, T-1$}{

  \BlankLine
  \tcp{Step 1: Client selection (coalition formation)}
  $S_t \gets \{i : w_i > c_{\min}\}$ \tcp*{exclude low-evidence clients}
  $W \gets \sum_{i \in S_t} w_i$\;
  $\tilde{w}_i \gets w_i / W$ for each $i \in S_t$\;
  Broadcast $\theta^t$ to all $i \in S_t$\;

  \BlankLine
  \tcp{Step 2: Local training}
  \For(\tcp*[f]{in parallel}){each $i \in S_t$}{
    $\theta_i \gets \theta^t$\;
    \For{$k = 0, \ldots, \tau - 1$}{
      Sample mini-batch $\xi \sim \cD_i$\;
      $\theta_i \gets \theta_i - \eta\,\nabla F_i(\theta_i; \xi)$\;
    }
    $\Delta_i^t \gets \theta_i - \theta^t$\;
  }

  \BlankLine
  \tcp{Step 3: Evidence-aware compression}
  \For{each $i \in S_t$}{
    Allocate bandwidth $R_i \propto \tilde{w}_i \norm{\Delta_i^t}$\;
    $k_i \gets \lfloor d \cdot R_i / R_{\text{total}} \rfloor$ \tcp*{coordinates to send}
    $\hat{\Delta}_i^t \gets \textsc{TopK}(\Delta_i^t, k_i)$ \tcp*{sparsification}
  }

  \BlankLine
  \tcp{Step 4: Evidence-weighted aggregation}
  $\theta^{t+1} \gets \theta^t + \sum_{i \in S_t} \tilde{w}_i\,\hat{\Delta}_i^t$\;

  \BlankLine
  \tcp{Step 5: Periodic variance re-estimation}
  \If{$t \bmod T_{\text{re-est}} = 0$}{
    Update $\hat{\sigma}_i^2$ from recent gradient statistics\;
    Recompute $w_i, \tilde{w}_i$\;
  }
}
\Return{$\theta^T$}\;
\end{algorithm}


%% ========================================================================
\section{Discussion}
\label{sec:discussion}
%% ========================================================================

\subsection{Variance Estimation in Practice}

The evidence weights $w_i = \sigma_{\min}^2 / \sigma_i^2$ require estimates
of client-level gradient variance. Two practical approaches:
\begin{enumerate}[(i)]
  \item \textbf{Pilot phase}: at initialization, each client computes
    multiple gradient estimates at the same $\theta^0$ and reports the
    empirical variance to the server. Cost: $O(M)$ extra gradient
    computations per client, where $M$ is the number of pilot samples.
  \item \textbf{Online estimation}: the server tracks the running variance
    of each client's updates $\Delta_i^t$ across rounds. After an initial
    burn-in period, these estimates stabilize and can be used for weighting.
    This requires no extra computation.
\end{enumerate}

Both approaches preserve privacy: the server only sees gradient statistics,
not raw data.

\subsection{Connections to the $\Omega$-Framework}

The FL application demonstrates the generality of the $\Omega$-framework's
core principles:
\begin{enumerate}[(1)]
  \item \textbf{Evidence weighting} ($w_i = V_{\min}/V_i$): the same
    weighting scheme improves convergence in both multi-agent RL and FL.
    The $\HM/\AM$ improvement factor is a universal consequence of optimal
    inverse-variance weighting.
  \item \textbf{Self-knowledge bottleneck}: the insight that ``you cannot
    communicate what you do not know'' applies equally to RL agents
    communicating policies and FL clients communicating gradients.
    Bandwidth cannot substitute for data quality.
  \item \textbf{Coalition formation}: client selection in FL is structurally
    identical to coalition formation in cooperative games. The evidence
    weight serves as the natural inclusion criterion in both settings.
  \item \textbf{Blessing of dimensionality}: sparsified gradient communication
    exploits the same concentration phenomena that make high-dimensional
    policy gradient tractable.
\end{enumerate}

\subsection{Limitations}

\begin{enumerate}[(i)]
  \item \textbf{Variance estimation overhead}: the pilot phase adds
    communication rounds. Online estimation introduces a burn-in period
    during which weights may be inaccurate.
  \item \textbf{Non-stationarity}: client data distributions may shift over
    time, requiring periodic re-estimation of $\sigma_i^2$.
  \item \textbf{Privacy}: reporting $\hat{\sigma}_i^2$ leaks information
    about local data heterogeneity. Differential privacy mechanisms can
    be applied to the variance estimates at the cost of noisier weights.
  \item \textbf{The evidence weight is not the only factor}: in practice,
    client reliability (uptime, network quality) and fairness considerations
    may override pure evidence-weighted selection.
\end{enumerate}


%% ========================================================================
\section{Conclusion}
\label{sec:conclusion}
%% ========================================================================

We have shown that the $\Omega$-framework's evidence-weighted multi-agent
policy gradient transfers precisely to federated learning. The core
mechanism --- inverse-variance weighting via $w_i = V_{\min}/V_i$ ---
produces a $\HM(V)/\AM(V)$ variance improvement that is maximally
beneficial precisely when data heterogeneity is worst. The self-knowledge
bottleneck explains why bandwidth upgrades cannot fix the non-IID problem:
clients with poor local data have fundamentally limited information about
the global optimum, regardless of communication capacity. Evidence-weighted
client selection provides a principled, automatic alternative to ad hoc
selection heuristics.

The FL application validates the $\Omega$-framework's claim to generality:
its components (evidence weighting, lossy communication, coalition
formation, blessing of dimensionality) are not artifacts of the multi-agent
RL setting but universal principles of distributed learning under
heterogeneous information quality.


%% ========================================================================
\bibliographystyle{plain}
\begin{thebibliography}{10}

\bibitem{mcmahan2017}
B.~McMahan, E.~Moore, D.~Ramage, S.~Hampson, and B.~A.~y~Arcas.
\newblock Communication-efficient learning of deep networks from decentralized
  data.
\newblock In \emph{Proceedings of AISTATS}, 2017.

\bibitem{karimireddy2020}
S.~P.~Karimireddy, S.~Kale, M.~Mohri, S.~Reddi, S.~Stich, and A.~T.~Suresh.
\newblock {SCAFFOLD}: Stochastic controlled averaging for federated learning.
\newblock In \emph{Proceedings of ICML}, 2020.

\bibitem{li2020fedprox}
T.~Li, A.~K.~Sahu, M.~Zaheer, M.~Sanjabi, A.~Talwalkar, and V.~Smith.
\newblock Federated optimization in heterogeneous networks.
\newblock In \emph{Proceedings of MLSys}, 2020.

\bibitem{giannou2022}
A.~Giannou, K.~Lotidis, P.~Mertikopoulos, and
  P.~Vlatakis-Gkaragkounis.
\newblock On the convergence of policy gradient methods to {N}ash equilibria in
  general stochastic games.
\newblock \emph{arXiv:2210.08857}, 2022.

\bibitem{ewpg2026}
E.~Shcherbinin.
\newblock Evidence-weighted policy gradient convergence in general stochastic
  games.
\newblock Working paper, 2026.

\bibitem{blessing2026}
E.~Shcherbinin.
\newblock The blessing of dimensionality for policy gradients: sparse recovery,
  evidence weighting, and the logarithmic cost of action space expansion.
\newblock Working paper, 2026.

\bibitem{cooperative2026}
E.~Shcherbinin.
\newblock Cooperative policy gradient with lossy communication in general
  stochastic games.
\newblock Working paper, 2026.

\bibitem{rothchild2020}
D.~Rothchild, A.~Panda, E.~Ullah, N.~Ivkin, I.~Stoica,
  V.~Braverman, J.~Gonzalez, and R.~Arber.
\newblock {FetchSGD}: Communication-efficient federated learning with
  sketching.
\newblock In \emph{Proceedings of ICML}, 2020.

\bibitem{li2019convergence}
X.~Li, K.~Huang, W.~Yang, S.~Wang, and Z.~Zhang.
\newblock On the convergence of {FedAvg} on non-{IID} data.
\newblock In \emph{Proceedings of ICLR}, 2020.

\end{thebibliography}

\end{document}
